\section{Problem description and two-stage scheduling model}
\setcounter{equation}{0}

All material considered here consists of wafers: product wafers and auxiliary (dummy) wafers. Products receive their prescribed recipe; auxiliary wafers fill production positions or occupy the chamber during cleaning and are reused after return. A chamber operation is either production or cleaning. Production means recipe processing in the $4\times1$ or $2\times2$ mode; cleaning uses only auxiliary wafers. One production or cleaning operation constitutes one chamber cycle. A production cycle specifically denotes a recipe-processing cycle and excludes cleaning. Loading, unloading, and rotation are separate actions; waiting is the interval between consecutive actions.

Given product counts and recipes at LP1--LP4, auxiliary-wafer inventory, and the initial equipment state, the scheduling task jointly determines wafer pairing, chamber assignment, chamber cycles, shared-resource orders, and event times. The model is organized by physical routes, wafer-to-cycle assignments, timing logic, production and cleaning rules, resource occupancy, and the two-stage objective.

Inputs comprise action durations, resource capacities, residency limits, the cleaning threshold, and auxiliary-wafer inventory. Stage one minimizes the time at which the last product returns to its LP. Stage two improves waiting and production cadence under a prescribed makespan bound. A feasible schedule specifies the wafers, resources, and start and end times of each action for performance evaluation and checks of precedence and resource conflicts.

The cleaning threshold satisfies $H\ge2$, and the shared inventory contains at least four auxiliary wafers. Action durations are deterministic within each solve. No buffers are assumed beyond LPs, aligners, load locks, rotating chambers, and vacuum storage.

Table~\ref{tab:symbol-en-a} summarizes sets, inputs, and capacity bounds; Table~\ref{tab:symbol-en-b} lists core decisions, derived indicators, and linearization auxiliaries. All events use start and end times $s_o,e_o$. For a fixed-duration physical action, $e_o=s_o+d_ob_o$ can be substituted directly. For a resource-holding interval spanning several actions, the endpoints are linked to the corresponding pickup and placement events, and its duration is not prescribed. Loading, unloading, rotation, cycle, and cleaning-stage times are indexed aliases of these physical events, not independent clocks.

Subscripts $w,i,p,c,q$ denote a product wafer, wafer-route task, PW pair, chamber, and candidate chamber cycle, respectively; $o,o'$ denote actions or resource-occupancy tasks, and $\alpha$ denotes a waiting item. In $a_{wpr},d_{pr}$, $r=0,1$ identifies a member position within a pair; in $y^r_{oo'}$, $r$ identifies a resource. The index $\ell\in\Gamma_o$ identifies a physical hand or slot. A loan request $j=(p,r)$ identifies one use of an auxiliary wafer, whereas $\tau$ identifies a reusable physical wafer. Cleaning segment $e$ and production task $v$ are derived from the cycle sequence and used to describe cleaning windows.

The labels $\zeta,\psi,\mathrm{cl}$ denote the two production recipes and cleaning. The labels load, unload, rot$\pm$, ref, return, and max denote loading, unloading, rotation before/after a cycle, a reference stage, product return, and a maximum. They are semantic labels rather than additional decision indices. All times are in seconds. A binary indicator is 1 when its stated choice holds and 0 otherwise. Derived quantities abbreviate expressions used in several constraints; their defining relationships introduce no independent decision freedom.

@@TABLE1@@

@@TABLE2@@

\subsection{Single- and dual-source physical routes}
\label{sec:model-path-en}

Let $\mathcal W$ be the product-wafer set and $\mathcal I$ the set of wafer-route tasks. Each product has one complete route; each loan of an auxiliary wafer has one supply-and-return route. One physical auxiliary wafer can serve several successive loan tasks.

Let $\mathcal S=\{A,V\}$ denote atmospheric and vacuum source regions. The source configuration is $\sigma\in\{0,1\}$ for single- and dual-source supply, respectively. LP1--LP4 specify product locations within the atmospheric region; each product returns to its original LP. Products satisfy $\chi_{wA}=1,\chi_{wV}=0$. Auxiliary task $i$ satisfies $\chi_{iA}=1-\sigma,\chi_{iV}=\sigma$. Source-membership parameters obey
@@eq:source-membership@@
Equation~\eqref{eq:source-membership-en} assigns a source region to every route task. Under single-source supply, products occupy LP1--LP4 and auxiliary wafers are stored beside LP4 in the atmospheric region. Both share ATR and load-lock interfaces, but only products require AL alignment. Under dual-source supply, products still originate at LP1--LP4, whereas auxiliary wafers originate in vacuum storage. Source configuration concerns the regions supplying the two wafer types, not the number of LPs.

In both configurations, products follow LP--ATR--AL--ATR--upper load lock--VTR--process chamber--VTR--lower load lock--ATR--original LP. Alignment precedes the load-lock pressure transition and execution of the prescribed recipe in the chamber.

Single-source auxiliary wafers travel from storage beside LP4 through ATR, the upper load lock, and VTR to the chamber, and return through VTR, the lower load lock, and ATR. A separately supplied auxiliary wafer bypasses AL. During an accompanying entry with a product, it remains on ATR while the product is aligned, and both enter the lock after product retrieval. Neither route aligns the auxiliary wafer.

Dual-source auxiliary wafers travel from vacuum storage through VTR to the chamber and return through VTR after use. Figures~\ref{fig:single-source-en} and \ref{fig:system-en-new} show both configurations. These routes specify the required physical stages; paired-transfer choices and their timing conditions are given in Section~\ref{sec:model-time-en}.

\subsection{Assignments of wafers, PW pairs, chambers, and cycles}
\label{sec:model-assignment-en}

Individual wafers are first organized into PW pairs, which are assigned to compatible chambers and chamber cycles. Production and cleaning use the same candidate-cycle set. Exposure counts and cleaning triggers are specified in Section~\ref{sec:model-process-en}.

\noindent\textbf{Wafer pairing and chamber selection.}
A wafer pair (PW pair) contains two products, one product and one auxiliary wafer, or two auxiliary wafers. Product LP $\ell_w$ and recipe $m_w$ are given; same-recipe products may be paired across LPs. The model chooses the members of each candidate pair $p\in\mathcal P$. Binary $a_{wpr}$ assigns product $w$ to member position $r\in\{0,1\}$, $d_{pr}$ indicates auxiliary occupancy, and $v_p$ activates the pair:
@@eq:joint-membership@@
Equation~\eqref{eq:joint-membership-en} assigns every product to exactly one pair and fills both positions of every active pair.

Let $\mathcal M=\{\zeta,\psi,\mathrm{cl}\}$ denote $4\times1$ production, $2\times2$ production, and cleaning. The set $\mathcal C(m)$ contains chambers compatible with operation type $m$. Binary variables $\mu_{pm}$ and $z_{pc}$ choose the pair's operation type and chamber and prohibit incompatible combinations:
@@eq:joint-chamber@@
Product recipe $m_w$ and source configuration $\sigma$ are inputs; pairing, chamber assignment, and operation order are decisions. Binary $B_p$ identifies a product-containing pair and $B_p^{\psi}$ a pair containing $2\times2$ products:
@@eq:pair-type@@

\noindent\textbf{Candidate-capacity bounds.}
Auxiliary-only pairs provide filling or cleaning; product-containing pairs cannot perform cleaning. Let $N=|\mathcal W|=N_{\zeta}+N_{\psi}$ be the number of products, excluding reusable auxiliary wafers. Here $N_{\zeta},N_{\psi}$ count products of the two recipes. The following conservative capacities reserve enough indices for joint pairing and cycle decisions:
@@eq:candidate-bounds@@
The bound $T$ limits the total number of production cycles across the tool. Each $4\times1$ product receives one exposure and each $2\times2$ product receives two; every production cycle contains at least one product. Thus the total number of production cycles is at most $N_{\zeta}+2N_{\psi}$. The threshold $H$ limits the production count in each chamber since its initially clean state or latest cleaning. Cleaning is optional at a valid boundary after 1 to $H-1$ production cycles and mandatory when the count reaches $H$. Initial empty cleaning and consecutive cleanings are prohibited, so each cleaning is preceded by at least one production cycle. Hence $J_{\max}=T$ safely bounds total cleanings. This bound retains all valid early-cleaning choices; it does not require $T$ cleanings.

The capacity $P_{\max}$ bounds the candidate pair set $\mathcal P=\{1,\ldots,P_{\max}\}$ for the entire lot. There are at most $N$ product-containing pairs because each contains a product and every product belongs to exactly one pair. Each production cycle contains two PW pairs and at least one product-containing pair, giving at most $T$ auxiliary-only support pairs. Each cleaning uses two auxiliary-only pairs, giving at most $2J_{\max}$ cleaning pairs. Their sum yields $P_{\max}$ in Eq.~\eqref{eq:candidate-bounds-en}; actual pair activation satisfies $\sum_{p\in\mathcal P}v_p\le P_{\max}$.

Each PW pair denotes one uninterrupted occupancy from loading to unloading. A product-containing $2\times2$ pair participates in two adjacent production cycles but uses one pair index. A physical auxiliary wafer can be borrowed by another pair after return, so candidate-pair capacity differs from physical auxiliary inventory. The bound $Q_{\max}=T+J_{\max}$ limits total production-plus-cleaning cycles across the tool and also supplies the candidate capacity for each chamber.

\noindent\textbf{Pair-to-cycle assignments.}
Let $\mathcal Q_c=\{1,\ldots,Q_{\max}\}$ be the candidate cycles of chamber $c$, with types $4\times1$, $2\times2$, and cleaning. The variable $K_{cqm}$ selects type, and $U_{cq}=\sum_mK_{cqm}\in\{0,1\}$ abbreviates cycle activation. The variable $X_{pcq}$ selects participating pairs, and $X^P_{pcq}$ identifies occupancy by a product-containing pair. Each active cycle contains two pairs with compatible type and chamber:
@@eq:cycle-capacity@@
The product-occupancy indicator is linked linearly by
@@eq:cycle-products@@
The indicators $B_p,B_p^{\psi},X^P_{pcq}$ and the later $A^{\psi}_{pcq},Z^{\psi}_{pcq}$ are logical auxiliaries determined by membership, type, and occupancy. Linear linking constraints are retained to avoid products of binary variables. Production cycles therefore contain a product pair, whereas cleaning cycles contain only auxiliary pairs. Prefix activation removes empty cycle indices without prescribing the order of operation types.

\subsection{Action timing and precedence logic}
\label{sec:model-time-en}

Given physical routes and pairing relationships, activation, durations, precedence, and synchronization describe route feasibility and technical action order. Orders chosen to resolve shared-resource conflicts are addressed in Section~\ref{sec:model-resource-en}.

\noindent\textbf{Action activation and duration.}
For candidate action $o\in\mathcal O$, let $b_o$ indicate activation, $s_o,e_o$ be start and end times, $d_o$ be its prescribed duration, and $M$ be a valid large constant. For a fixed-duration physical action, the end time follows from the start time and activation:
@@eq:operation-time@@
Inactive actions have zero times; active actions last $d_o$. This form covers robot transfers, alignment, pumping, venting, rotation, recipe processing, auxiliary-wafer handling, and cleaning.

\noindent\textbf{Route-specific precedence.}
Let $\mathcal A^{\rm int}$ contain common chamber-internal precedence arcs, and let $\mathcal A_{is}(\gamma)$ contain the supply-and-return arcs for route task $i$, source region $s$, and transfer choices $\gamma$. The latter collectively denote paired transfers and single-source accompanying entry. The complete arc set is
@@eq:route-union@@
Equation~\eqref{eq:route-union-en} activates the route appropriate to each source and transfer choice and connects product routes, auxiliary supply/return routes, and chamber action chains.

For each required precedence arc $(o,o')\in\mathcal A$,
@@eq:precedence@@
A successor can execute only if its predecessor is active. When both are active, the successor starts after predecessor completion and the minimum interval $\delta_{oo'}$. These constraints connect wafer routes and loading--operation--unloading chains.

\noindent\textbf{Stage links and paired-action synchronization.}
An individual wafer's stage must coincide with the stage selected for its batch, chain position, or cleaning task. For selection indicator $x$ and reference interval $[s^{\rm ref},e^{\rm ref}]$,
@@eq:stage-link@@
When $x=1$, the individual action and reference stage coincide; when $x=0$, this candidate link is inactive. The link connects pair positions, product routes, auxiliary requests, and chamber operation windows in one time model.

Let $w_1,w_2$ be the two wafers in pair $p$. The variable $\gamma_{pa}$ selects a paired transfer on an eligible atomic transfer leg $a$ at a common location; $s_{wa},e_{wa}$ denote wafer $w$'s start and end times on that leg. Paired transfer requires synchronization:
@@eq:double-transfer@@
When $\gamma_{pa}=1$, the robot transfers the complete pair in one nonpreemptive action. Otherwise, the two wafer actions still use the same robot and are ordered by Eq.~\eqref{eq:resource-order-en}. Actuator exclusivity counts coordinated actions: the two synchronized wafer events form one action and are not ordered against each other, while their common action remains exclusive with other actions. Each wafer retains its own hand-occupancy interval, so synchronization does not merge hand capacity.

\noindent\textbf{Entry logic for a single-source product--auxiliary pair.}
Let $w_1$ be the product and $w_2$ the auxiliary wafer. The variable $\gamma_{p,\mathrm{in}}$ chooses accompanying entry after both are collected. This route link does not require simultaneous pickup from different storage positions. The event labels $\mathrm{pick,ALput,ALpick,LLput}$ denote source pickup, placement on AL, retrieval from AL, and placement into the upper load lock. The indicators $\gamma^{\rm early}_p$ and $\gamma^{\rm late}_p$ choose direct auxiliary entry before product alignment or auxiliary pickup after product placement on AL. Exactly one of these and accompanying entry is selected:
@@eq:single-source-accompany@@
For accompanying entry, ATR holds the auxiliary wafer before the product is placed on AL, and both enter the load lock together after alignment. The auxiliary wafer occupies one ATR hand from source pickup through completion of load-lock placement. The product occupies AL rather than an ATR hand during alignment and resumes hand occupancy after retrieval. Separate auxiliary supply directly links its pickup and load-lock placement. Neither option activates an AL-alignment event for the auxiliary wafer.

\subsection{Production modes and cleaning}
\label{sec:model-process-en}

Section~\ref{sec:model-assignment-en} determines which pairs occupy each cycle. This section specifies exposure counts, retention and departure between cycles, rotation and transfer order, and cleaning boundaries. All times use the physical-action representation in Section~\ref{sec:model-time-en}.

\noindent\textbf{Production modes and continuous occupancy.}
A regular all-product $4\times1$ batch contains two product PW pairs. Four products receive one synchronous exposure and leave in last-in--first-out order; auxiliary wafers fill shortages. In $2\times2$, each product-containing pair occupies two adjacent production cycles in the same chamber, giving every product two exposures. For three product pairs, a leading auxiliary pair D is loaded, the chamber rotates, and P1 is loaded for production on D+P1. Rotation and chain exchanges then produce P1+P2, P2+P3, and P3+D. P1, P2, and P3 each receive two exposures; the trailing auxiliary pair preserves four-position occupancy. Figures~\ref{fig:chain-en} and \ref{fig:parallel-en} show the sequences.

@@FIGCHAIN@@

@@FIGPARALLEL@@

Binary $A_{pcq},Z_{pcq}$ identify the first and last cycles occupied by pair $p$ in chamber $c$. With $X_{pc0}=X_{pc,Q_{\max}+1}=0$,
@@eq:cycle-exposures@@
Here $A^{\psi}_{pcq}=A_{pcq}\land B_p^{\psi}$ and $Z^{\psi}_{pcq}=Z_{pcq}\land B_p^{\psi}$ use the three-inequality linearization in Eq.~\eqref{eq:cycle-products-en}. Each active pair enters and exits once, ensuring uninterrupted occupancy. Total occupancy $v_p+B_p^{\psi}$ gives one production cycle to a $4\times1$ product pair and two adjacent production cycles in one chamber to a $2\times2$ product pair. Each $2\times2$ cycle contains at most one product pair receiving its first exposure and one receiving its last, forming an alternating exchange chain. An auxiliary-only support pair occupies one production cycle, and a cleaning pair occupies one cleaning cycle.

\noindent\textbf{Cycle durations and timing links.}
Let $S_{cq},E_{cq}$ be cycle start and end times, $d_m$ the duration of operation type $m$, and $\overline T$ a valid time horizon:
@@eq:cycle-time@@
Cycle durations are fixed; waiting occurs between actions. Default durations are 180~s for $4\times1$ production, 90~s for $2\times2$ production, and 500~s for cleaning. Transfer and rotation durations are counted through their separate events.

\noindent\textbf{Front/back occupancy and rotating-chamber transfers.}
Binary $f_{pcq}$ selects the pair at the robot-accessible front side during a cycle. The derived binary indicators $\kappa_{cq},\upsilon_{cq}$ denote retention of one pair from the preceding cycle and complete departure after the current cycle. They follow from occupancy changes rather than independent choices:
@@eq:cycle-side@@
A newly loaded pair occupies the front side. A retained pair rotates to the back and must have occupied the front in the preceding cycle. At a new $2\times2$ chain head, the product pair is loaded last into the front. Pair loading and unloading times are $s_p^{\rm load},e_p^{\rm load},s_p^{\rm unload},e_p^{\rm unload}$. Rotation times before and after cycle $(c,q)$ are $s_{cq}^{\rm rot\pm},e_{cq}^{\rm rot\pm}$. These are aliases of physical events. Rotations $\mathrm{rot-}$ and $\mathrm{rot+}$ are activated by $U_{cq}$ and $\upsilon_{cq}$, respectively, and last $r_c$. Loading links are
@@eq:cycle-load@@
When both pairs leave, post-cycle rotation starts after the chamber operation ends. Unloading after the last exposure obeys
@@eq:cycle-unload@@
If the next cycle retains a pair, its new pair waits until the replaced back pair has been unloaded; otherwise, the next loading waits for both previous pairs to leave. Indicator constraints or valid big-$M$ inequalities linearize these conditions, together with VTR exclusivity. Waiting remains explicit rather than being absorbed into rotation or transfer durations.

\noindent\textbf{Cleaning triggers and action chains.}
Integer $n_{cq}$ counts production cycles since the latest cleaning at completion of cycle $q$. With terminal candidate activation $U_{c,Q_{\max}+1}=0$,
@@eq:clean-counter@@
Each production cycle increments the count by one, and cleaning resets it to zero. Cleaning is mandatory at $H$. After 1 to $H-1$ production cycles, it may be selected at a valid boundary satisfying full occupancy, uninterrupted exposures, and resource constraints, with no restriction on the next production mode. Thus successive cleanings are separated by 1 to $H$ production cycles. A final segment below $H$ need not be followed by cleaning; reaching $H$ requires terminal cleaning. Equations~\eqref{eq:cycle-capacity-en} and \eqref{eq:cycle-exposures-en} keep the two exposures of each $2\times2$ product pair adjacent, so cleaning cannot split them.

Cleaning segments $e$ follow from cleaning boundaries in the cycle-type sequence. Let $\mathcal V_c$ contain production tasks, each being a $4\times1$ batch or a $2\times2$ production cycle. The indicator $a_{vce}$ assigns task $v$ to segment $e$ of chamber $c$, and $q^{\rm cl}_{ce}$ indicates cleaning after that segment. These indicators are linked to the cycle structure rather than providing independent cleaning-count or segmentation decisions. Internal segments end with cleaning; terminal cleaning follows Eq.~\eqref{eq:clean-counter-en}.

Task $v$ uses event times $s_v,e_v$. The complete cleaning task following segment $e$ has endpoints $s^{\rm cl}_{ce},e^{\rm cl}_{ce}$; its stage times are $s^{\rm cl,\nu}_{ce},e^{\rm cl,\nu}_{ce}$ for $\nu\in\{FL,BL,P,BU,FU\}$, denoting front loading, back loading, cleaning treatment, back unloading, and front unloading. These aliases refer to physical events; the complete task spans front-loading start through front-unloading completion. The timing relationships are
@@eq:clean-time@@
Segment $e$ starts after the preceding cleaning task ends and ends before the next cleaning begins. A complete cleaning task includes loading auxiliary wafers on both sides, rotation, cleaning treatment, and unloading in reverse order. Its durations and resource occupancy constrain VTR, the chamber, and all four auxiliary wafers.

Once a solution is given, $4\times1$ batch sides, $2\times2$ product-pair chain positions, and cleaning segments can be recovered from $X,A,Z,f$ and cycle types without additional independent position variables. Fixed-layout experiments additionally fix these positions and service orders and permit filling only in the final short batch: if later same-type batches exist, the current $4\times1$ batch contains two product pairs. The joint model allows auxiliary filling in any valid production cycle.

\subsection{Auxiliary-wafer reuse and robot/module resources}
\label{sec:model-resource-en}

Wafer routes and chamber rules must also respect physical-resource availability. Each auxiliary loan is first assigned a physical identity, followed by physical-hand and module-slot assignments. Resource orders and residency limits coordinate their occupancy.

\noindent\textbf{Auxiliary identities and reuse.}
Each auxiliary member position defines a loan request $j=(p,r)$, activated directly by $d_{pr}$ for product--auxiliary, auxiliary-only filling, or cleaning pairs. Membership and cycle choices link its supply and return events. Physical auxiliary occupancy extends from storage pickup through completed return.

All chambers access one physical auxiliary inventory. The set $\mathcal T_c$ contains identities accessible to chamber $c$; loans of the same identity must be exclusive across chambers as well. Each active request selects one physical wafer:
@@eq:dummy-assign@@
The assignment $g_{j\tau}$ also marks occupancy of the borrowed physical identity. Shared-resource constraints order successive loans, permitting reuse only after unloading and completed return to storage. Identity choices determine which cross-chamber loans compete: even with identical product pairing, chamber assignment, and cycle structure, different identities can change concurrent production intervals. Section~\ref{sec:joint-evidence-en} quantifies this effect with product structure controlled.

\noindent\textbf{Robot hands and module capacities.}
ATR has two hands, VTR has four, and each rotating chamber has four wafer positions. At any instant, each ATR or VTR executes at most one coordinated transfer, and each chamber executes at most one unified loading, unloading, rotation, or chamber-operation stage. Their actuators are therefore exclusive action resources. Individual hand assignments express robot carrying capacity, and pair occupancy expresses chamber capacity; actuator exclusivity does not represent wafer count. AL is also exclusive. The upper and lower load locks contain slots U1/U2 and L1/L2, respectively, and every physical auxiliary wafer is a reusable identity. Each occupancy interval selects one allowed unit, whether a physical hand, an auxiliary identity, or an exclusive actuator. For occupancy task $o$,
@@eq:slot-select@@
Binary $b_o$ activates the task, $h_{o\ell}$ selects unit $\ell$, and $\Gamma_o$ is the set of allowed units. Inactive tasks select no unit. Synchronized wafers retain separate hand occupancies, with synchronization enforced by Eq.~\eqref{eq:double-transfer-en}. Full chamber occupancy follows from pair assignments. Equation~\eqref{eq:resource-order-en} prevents overlapping occupancy of any physical unit. If $\Gamma_o$ is a singleton, $h_{o\ell}=b_o$ can be substituted directly. Auxiliary-identity occupancy uses $g_{j\tau}$ without a duplicate assignment decision.

Upper-slot occupancy extends from ATR placement start through VTR pickup completion plus reset time. Lower-slot occupancy extends from VTR placement start through ATR pickup completion plus reset time. Reset limits reuse of that slot and may overlap subsequent robot travel.

VTR has dirty hands $D_1,D_2$ and clean hands $C_1,C_2$. Product holding from upper-lock pickup to chamber placement permits $\Gamma_o=\{D_1,D_2\}$; holding after all required exposures, from chamber pickup to lower-lock placement, permits $\Gamma_o=\{C_1,C_2\}$. Auxiliary input and return permit all four hands. Each holding interval lasts from pickup start through placement completion and retains the same physical hand selected by Eq.~\eqref{eq:slot-select-en}. Same-hand intervals are exclusive with zero hand-switching gap; robot travel is separately included through actuator constraints. Each hand class therefore holds at most two wafers, with auxiliary wafers competing for both classes. A $2\times2$ product remains inside the chamber between its two exposures and leaves on a clean hand after the second. Paired exchange permits dirty hands to hold an incoming pair while clean hands remove a completed pair, provided every auxiliary wafer is also assigned a free physical hand.

\noindent\textbf{Resource exclusivity and location-dependent travel.}
Resource orders are either prescribed or selected. Precedence arcs and predefined service chains impose the former; order variables resolve other conflicts. Fixed-layout pure-$4\times1$ experiments prescribe aligner and load-lock service orders and, when the number of candidate batches does not exceed the cleaning threshold, a VTR batch-transfer chain. These restrictions apply only to fixed-layout experiments; the joint model chooses the remaining resource orders.

For tasks $o,o'$ competing for unit $\ell$ of resource $r$,
@@eq:resource-order@@
When both tasks select the unit, binary $y^r_{oo'}$ chooses their order. The parameter $\rho^r_{oo'}$ includes repositioning or reset time. The constraint covers ATR, AL, load locks, VTR, chambers, and physical auxiliary identities, preventing conflicting actuator actions as well as premature reuse of slots or wafers.

Travel times are specified by location pairs. For actions whose chamber is chosen by $z_{pc}$, the relevant travel gap is activated by the selected chamber combination rather than a predetermined location. For example, when active exclusive chamber actions $o,o'$ serve pairs $p,p'$,
@@eq:location-travel@@
The parameter $\rho^{\rm VTR}_{cc'}$ is the travel time from chamber $c$ to $c'$. Reverse order is handled symmetrically; fixed-location actions use Eq.~\eqref{eq:resource-order-en} directly.

\noindent\textbf{Residency limits and timing feasibility.}
For the events at which wafer $w$ enters and leaves constrained equipment region $r$,
@@eq:residence@@
This bounds physical residency in a load lock, chamber, or constrained transport segment. Load-lock slot occupancy ends after the wafer is picked up from the opposite side; subsequent pressure reset may overlap robot travel, while the next incompatible load-lock action must await reset completion.

The precedence constraints in Eq.~\eqref{eq:precedence-en} and resource-order constraints in Eq.~\eqref{eq:resource-order-en} jointly prevent modeled deadlocks. Construct a directed graph whose nodes are active actions and whose arcs are technical precedence or selected resource orders. A circular wait among robots, load locks, chambers, or auxiliary wafers creates a directed cycle. Summing its time inequalities accumulates a positive duration or switching increment for actual occupancy actions, requiring an action to start strictly later than itself. Such an order is infeasible. Together with explicit load-lock slots, full chamber occupancy, and return-before-reuse requirements, time feasibility excludes circular waiting among the modeled actions and permits wafers to complete their routes.

\subsection{Makespan-first two-stage optimization}
\label{sec:model-objective-en}

The preceding constraints define the joint feasible region $\mathcal F$, with $\xi$ the vector of retained variables. Stage one optimizes product-return makespan. Stage two improves waiting and production cadence under the first-stage makespan bound and principal discrete structure.

\noindent\textbf{Stage one: minimize product-return makespan.}
The return-action end time $e_{w,\mathrm{return}}$ directly denotes product $w$'s return to its original LP:
@@eq:primary-objective@@
Makespan is measured by the last product's completion of its entire route. Cleaning between production segments delays subsequent products through resource and precedence constraints and thus affects makespan. Terminal cleaning after the last production remains in the action sequence but its completion is not a direct lower bound on $C_{\max}$. It may still delay returns from another chamber through the shared VTR, so indirect effects require the complete schedule. WPH uses product-return $C_{\max}$ as its denominator. The joint feasible region includes pairing, chamber assignment, cycle occupancy, and complete physical timing; fixed-layout experiments add their layout and service-order restrictions.

\noindent\textbf{Stage two: restrict the refinement domain.}
Let $\widehat\xi$ be the best first-stage incumbent and $\widehat C$ its product makespan. The symbol $\xi_d$ denotes component $d$, and hatted quantities are constants in stage two. The index set $\mathcal D_{\rm fix}$ covers fixed pairing, chamber assignments, cycle occupancy, auxiliary identities, hand assignments, accompanying-entry choices, and resource orders. With nonnegative makespan tolerance $\varepsilon_C$, impose
@@eq:secondary-domain@@
These constraints retain principal discrete structure and the makespan bound, while same-location atomic-transfer synchronization and continuous times remain adjustable. Stage two can start from a first-stage incumbent even without a proof of optimality.

Let $\mathcal Q^w$ index the waiting items and $\omega_\alpha$ be the nonnegative interval remaining after subtracting necessary action or travel time between the corresponding events. Total waiting enters the objective directly as $\sum_\alpha\omega_\alpha$. The experimental label $L_{\rm sum}$ is only a summary statistic for this sum, not an independent variable. The quantity $L_{\max}$ bounds the largest individual wait. Simultaneous waits of different wafers are counted separately. Each $\omega_\alpha$ is an event-time difference expression rather than an independent waiting decision:
@@eq:secondary-wait@@

Cadence uses cycle start times $S_{cq}$ directly. The set $\mathcal I_h$ contains the central indices $(c,q)$ of three consecutive production cycles within the same chamber and $2\times2$ chain. Both boundaries retain a pair, so $\kappa_{cq}=\kappa_{c,q+1}=1$; differences do not cross cleaning or chain boundaries. Fixing cycle structure in stage two fixes this set. The difference between successive start-to-start intervals is $S_{c,q+1}-2S_{cq}+S_{c,q-1}$. Nonnegative $v_{cq}$ bounds its absolute value linearly without a separate segment clock or difference variable:
@@eq:secondary-rhythm@@

When soft waiting limits $\bar\omega_\alpha$ are enabled, nonnegative $\eta_\alpha$ bounds excess waiting:
@@eq:secondary-softwait@@

With nonnegative weights $\lambda_1,\ldots,\lambda_4$, the second-stage objective is
@@eq:secondary-objective@@
The terms penalize total waiting, maximum individual waiting, variation in adjacent production-cycle start intervals, and soft-limit excess. The experiments use $(\lambda_1,\lambda_2,\lambda_3)=(1,10,25)$ and $\varepsilon_C=10^{-6}$. Disabling soft waiting limits omits their variables and penalty.

At an optimal second-stage solution with strictly positive corresponding weights, the auxiliary bounds are tight: $L_{\max}=\max_\alpha\omega_\alpha$, $v_{cq}=|S_{c,q+1}-2S_{cq}+S_{c,q-1}|$, and $\eta_\alpha=\max(0,\omega_\alpha-\bar\omega_\alpha)$. Time-limited feasible solutions may retain slack, so reported metrics are recomputed from actual event times. Cadence variation is zero for three equally spaced production starts. Soft waiting limits allow penalized excess; residency limits remain hard constraints.
